% Tabla de objetivos quitada de Conclusiones (Síntesis de resultados), guardada para la presentación.
\begin{table}[htbp]
    \centering
    \small
    \begin{tabular}{p{4.2cm}p{2.6cm}p{7.2cm}}
        \toprule
        Objetivo específico & Secciones & Resultado principal \\
        \midrule
        1. Caracterizar el costo marginal del SEN & 3.1, \ref{sec:verificacion_heterocedasticidad} & Ciclo diario marcado por la generación solar, heterocedasticidad horaria y colas pesadas del error \\
        2. Diseñar las estrategias P+C y S+C & 3.4.1 & Ambas implementadas y evaluadas en las ocho barras \\
        3. Comparar LayerNorm y DyT & \ref{sec:ln_vs_dyt} & Sin preferencia: la diferencia es del orden de la variabilidad entre semillas \\
        4. Stacking Optimization en P+C & 3.5, \ref{sec:stacking_resultados} & Supera a Prophet en las ocho barras, pero no a S+C \\
        5. Comparar P+C y S+C & \ref{sec:pc_vs_sc}, \ref{sec:h24_conclusiones} & S+C es más preciso en 14 de 16 combinaciones a una hora (9 significativas tras Holm-Bonferroni); a 24 horas P+C es algo más preciso \\
        6. Conformal Prediction y LW-ACP & \ref{sec:lwacp_comparativa}, \ref{sec:lwacp_robusto} & LW-ACP obtiene el menor Winkler al 90\%; su variante robusta supera a ACI en las ocho barras \\
        7. Variables hídricas en Puerto Montt & \ref{sec:pmontt_desempeno} & No mejoran el desempeño respecto de P+C y S+C \\
        \bottomrule
    \end{tabular}
    \vspace{0.2cm}
    \caption{Objetivos específicos del trabajo, secciones donde se abordan y resultado principal de cada uno.}
    \label{tab:objetivos}
\end{table}
